\chapter*{Lời mở đầu}
\addcontentsline{toc}{chapter}{Lời mở đầu}

Sự phát triển của văn minh nhân loại cùng với tốc độ đô thị hóa và công nghiệp hóa nhanh chóng đã và đang từng ngày cải tiến cuộc sống của con người trên nhiều khía cạnh khác nhau, từ kinh tế, công nghệ và giáo dục. Tuy vậy, ô nhiễm môi trường đồng thời cũng là một hệ quả to lớn và đáng quan ngại của xu hướng phát triển trên, bất chấp mọi nỗ lực của toàn nhân loại trong việc ngăn chặn và giảm thiểu tác động xấu đến đời sống và sức khỏe con người. Một trong những nguyên nhân tác động tiêu cực nhất tới nhân loại chính là ô nhiễm không khí. Tình trạng này không chỉ gây ra những bất tiện trong sinh hoạt hằng ngày mà còn tiềm ẩn nhiều nguy cơ đe dọa trực tiếp đến sức khỏe và sự phát triển bền vững của xã hội.

Ô nhiễm không khí đặt ra thách thức đáng kể đối với sức khỏe cộng đồng tại Việt Nam, đặc biệt là ở các trung tâm đô thị như Hà Nội. Hàng năm tại Việt Nam, rất nhiều bệnh nhân tử vong do các bệnh liên quan đến ô nhiễm không khí, bao gồm bệnh tim, đột quỵ, ung thư phổi, bệnh phổi tắc nghẽn mãn tính và viêm phổi. Những nguyên nhân chính gây ô nhiễm không khí tại Việt Nam là quá trình đô thị hóa nhanh chóng và công nghiệp hóa, dẫn đến lượng khí thải từ giao thông, hoạt động công nghiệp và sử dụng năng lượng dân dụng tăng lên \cite{reason}. Bên cạnh đó, mật độ phương tiện giao thông dày đặc, khí thải công nghiệp và các hoạt động xây dựng của thành phố cũng góp phần làm trầm trọng thêm vấn đề, dẫn đến chất lượng không khí thường xuyên vượt quá giới hạn an toàn \cite{HIEN2020134635, Ly_2018}. 

Dù vậy, việc tiếp cận với các thông tin liên quan đến chất lượng không khí xung quanh khu vực sinh sống vẫn còn hạn chế đối với người dân. Đây là một thách thức lớn đặt ra nếu muốn triệt để ngăn chặn tình trạng ô nhiễm không khí kéo dài và có những biện pháp phù hợp và kịp thời để đẩy lùi những tác động tiêu cực. Nhu cầu giám sát chất lượng không khí và cung cấp thông tin cảnh báo kịp thời đến người dân đã trở thành động lực chính cho nghiên cứu này của chúng tôi. Trong dự án này, chúng tôi xây dựng hệ thống \textbf{HNAQI} - một hệ thống có khả năng theo dõi định kỳ Chỉ số Chất lượng Không khí (Air Quality Index) tại Hà Nội, kết hợp với cơ chế cảnh báo dựa trên các ngưỡng được định nghĩa rõ ràng. Hệ thống này không chỉ giúp người dùng cập nhật liên tục tình hình không khí trên toàn Hà Nội mà còn hỗ trợ đưa ra các quyết định phù hợp về hoạt động ngoài trời và kế hoạch tương lai gần, từ đó góp phần giảm thiểu các rủi ro không mong muốn, ví dụ như sức khỏe và các dự định ở thời điểm đó.

Quá trình xây dựng hệ thống \textbf{HNAQI} giám sát, cảnh báo và dự đoán chất lượng không khí tại Hà Nội (từ đây gọi tắt là \textbf{Hệ thống HNAQI}) sử dụng đầu vào là dữ liệu từ nhiều nguồn dữ liệu về chất lượng không khí khác nhau và có thể truy cập công khai. Trải qua quá trình truy xuất, chuyển đổi và lưu trữ vào cơ sở dữ liệu, các dữ liệu này sẽ được sử dụng trong các module cảnh báo và dự đoán phía sau. Đầu ra chính là một hệ thống có khả năng giám sát định kỳ và đưa ra cánh báo về chất lượng không khí quá khứ và hiện tại, đồng thời cung cấp dự báo tương lai gần theo mức giờ, mức ngày để giúp người dùng tiện theo dõi và sắp xếp lịch trình sao cho phù hợp.

Báo cáo này bao gồm ?? chương và được trình bày theo thứ tự sau:



